\subsection{Filters and masures}
This section aims to recall what masures are. As stated in the introduction, the reader only interested in the completion of Iwahori-Hecke algebras can skip this section and go directly to Section~4.

Masures were first introduced for symmetrizable split Kac--Moody groups over a valued field whose residue field contains $\jepPubBbb{C}$ by Gaussent and Rousseau~\cite{jep88refGR08}. Later, Rousseau axiomatized this construction in~\cite{jep88refRou11}, then generalized it with further developments to almost-split Kac--Moody groups over non-Archimedean local fields in~\cite{jep88refRou16,jep88refRou17}. For the reader familiar with this work, let us mention that we consider here semi-discrete masures which are thick of finite thickness.
\subsubsection{Filters, sectors and rays}
\begin{definition}\label{jep88localDefinition2_2}
A filter on a set $E$ is a non-empty set $F$ of non-empty subsets of $E$ that satisfies the following conditions:
\begin{itemize}
\item for any subsets $S_1,S_2$ of $E$ that both belong to $F$, then $S_1\cap S_2$ belongs to $F$;
\item for any subsets $S_1,S_2$ of $E$ with $S_1$ in $F$ and $S_1\subset S_2$, then $S_2$ belongs to $F$.
\end{itemize}
\end{definition}
Given a filter $F$ on a set $E$ and a subset $E'$ of $E$, we say that $F$ contains $E'$ if every element of $F$ contains $E'$. If $E'$ is non-empty, then the set $F_{E'}$ of all subsets of $E$ containing $E'$ is a filter on $E$ called the filter associated to $E'$. By language abuse and to ease notations, we will sometimes write that $E'$ is a filter, by identification of $E'$ with $F_{E'}$.

Now, let $F$ be a filter on a finite-dimensional real affine space $E$. We define its closure $\overline{F}$ as the filter of all subsets of $E$ that contain the closure of some arbitrary element of $F$, and its convex hull $\operatorname{conv}(F)$ as the filter of all subsets of $E$ that contain the convex hull of some arbitrary element of $F$. Said differently, we have:
\[
\overline{F}:=\{S\subset E\mid\exists S'\in F,\ \overline{S'}\subset S\}\quad\text{and}\quad\operatorname{conv}(F):=\{S\subset E\mid\exists S'\in F,\ \operatorname{conv}(S')\subset S\}.
\]
Given two filters $F_1$ and $F_2$ on the same set $E$, we say that $F_1$ is contained in $F_2$ iff any subset of $E$ contained $F_2$ is in $F_1$. Similarly, we say that the filter $F_1$ is contained in a subset $\Omega$ of $E$ iff any subset of $E$ contained in $\Omega$ is in $F_1$.

Let $\Omega$ be a subset of $\jepPubBbb{A}$ containing an element $x$ in its closure $\overline{\Omega}$ of $\Omega$. The germ of $\Omega$ in $x$ is defined as the filter $\operatorname{germ}_x(\Omega)$ of all subsets of $\jepPubBbb{A}$ containing some neighborhood of $x$ in $\Omega$. A sector in $\jepPubBbb{A}$ is a translate $\jepPubFrak{s}:=x+C^v$ of a vectorial chamber $C^v=\pm w\cdot C_f^v$ (with $w\in W^v$) by an element $x\in\jepPubBbb{A}$. The point $x$ is called the base point of the sector $\jepPubFrak{s}$ and the chamber $C^v$ is called its direction. One can easily check that the intersection of two sectors having the same direction is a sector with the same direction.

Given a sector $\jepPubFrak{s}:=x+C^v$ as above, the sector-germ of $\jepPubFrak{s}$ is the filter $\jepPubFrak{S}$ of all subsets of $\jepPubBbb{A}$ containing an $\jepPubBbb{A}$-translate of $\jepPubFrak{s}$. Note that it only depends on the direction $C^v$ of $\jepPubFrak{s}$. In particular, we denote by $\pm\infty$ the sector-germ of $\pm C_f^v$.

Finally, let $\delta$ be a ray with base point $x$ and let $y\neq x$ be another point on $\delta$ (which amounts to say that $\delta$ contains the interval $]x,y]=[x,y]\smallsetminus\{x\}$, as well as the interval $[x,y]$). We say that $\delta$ is pre-ordered (resp. generic) if either $y\leqslant x$ or $x\leqslant y$ (resp. if $y-x\in\pm\mathring{\jepPubScript{T}}$, where $\pm\mathring{\jepPubScript{T}}$ denotes the interior of the cone $\pm\jepPubScript{T}$).
